\section*{अध्याय \ref{ch:g12:subduction-magmatism} --- क्षेपण क्षेत्रों पर मैग्मा-क्रिया}

\begin{solution}{exo:g12:subduction-magmatism:1}
(a) \omterm{def:g12:subduction-magmatism:arc}{ज्वालामुखीय चाप}; (b) \omterm{def:g12:subduction-magmatism:hydrous-melting}{जलीय गलन} (गालक-प्रेरित गलन); (c) \omterm{def:g12:subduction-magmatism:granodiorite}{ग्रैनोडायोराइट};
(d) \omterm{def:g12:subduction-magmatism:batholith}{बैथोलिथ}; (e) \omterm{def:g12:subduction-magmatism:continental-growth}{महाद्वीपीय वृद्धि}।
\end{solution}

\begin{solution}{exo:g12:subduction-magmatism:2}
(1) कम से कम $24/0.25 = \qty{96}{km/h}$, यानी $\num{96000}/3600 \approx
\qty{27}{m/s}$। (2) विस्फोट कम से कम $480/96 = 5$ गुना तेज़ था।
\end{solution}

\begin{solution}{exo:g12:subduction-magmatism:3}
(1) $838/1214 \approx 0.690$, लगभग \qty{69}{\percent}। (2) ऊपर चढ़ने वाली
प्लेट की ओर, उस प्लेट की ओर जो नहीं डूबती। (3) खाई और
अग्र-रेखा के बीच \omterm{def:g11:subduction-zones:slab}{स्लैब} अभी पर्याप्त गहरी नहीं है: वह उस गहराई तक नहीं पहुँची है,
लगभग सौ किलोमीटर, जहाँ उसका जल ऐसे प्रावार में मुक्त होता है जो भीगने पर
पिघलने लायक गरम हो।
\end{solution}

\begin{solution}{exo:g12:subduction-magmatism:4}
(1) $45/(45 + 15) = 0.75$: फ़ेल्डस्पार का \qty{75}{\percent}
प्लेजियोक्लेज़ है। (2) हाँ: वह स्थूलकणी है, उसमें
\qty{20}{\percent} से अधिक क्वार्ट्ज़ है, और उसका प्लेजियोक्लेज़ अंश 65 और
\qty{90}{\percent} के बीच है।
\end{solution}

\begin{solution}{exo:g12:subduction-magmatism:5}
चाप-मैग्मा सिलिका में अधिक समृद्ध होते हैं, इसलिए वे अधिक श्यान होते हैं; और वे
उस जल को ढोते हैं जिसने उन्हें बनाया, इसलिए वे गैस से समृद्ध होते हैं। ऊँची श्यानता और
ऊँची गैस-मात्रा दोनों विस्फोटक उद्गारों को बढ़ावा देती हैं।
\end{solution}

\begin{solution}{exo:g12:subduction-magmatism:6}
(1) $H = 150 \times \tan\ang{35} \approx 150 \times 0.700 \approx
\qty{105}{km}$। (2) यह संसार का औसत ही है।
\end{solution}

\begin{solution}{exo:g12:subduction-magmatism:7}
$d = H/\tan\theta$: \ang{20} पर $105/\tan\ang{20} \approx 105/0.364 \approx
\qty{288}{km}$, और \ang{45} पर $105/\tan\ang{45} = \qty{105}{km}$।
\omterm{def:g11:subduction-zones:slab}{स्लैब} जितनी मंद, चाप खाई से उतनी ही दूर खड़ा होता है।
\end{solution}

\begin{solution}{exo:g12:subduction-magmatism:8}
(1) $T = -5.104 \times 3.16^2 + 132.9 \times 3.16 + 1120.7 \approx -51.0 +
420.0 + 1120.7 \approx \qty{1490}{\celsius}$। (2) लगभग $1490 - 1000 =
\qty{490}{\celsius}$। (3) शुष्क, नहीं: \qty{1300}{\celsius},
\qty{1490}{\celsius} से नीचे है। भीगी, हाँ: वह \qty{1000}{\celsius} से ऊपर है।
\end{solution}

\begin{solution}{exo:g12:subduction-magmatism:9}
(1) $2.5/2.7 \approx 0.93$ वर्ष, लगभग 340 दिन। (2) हिमस्खलन ने पहले से मौजूद
पुरानी चट्टान खिसकाई; चाप प्रावार से नई चट्टान बनाते हैं। केवल
आयतनों की तुलना की गई है।
\end{solution}

\begin{solution}{exo:g12:subduction-magmatism:10}
(1) $2.7 \times \num{1000000} = \qty{2.7e6}{km^3}$। (2) $\num{9e9}/2.7
\approx \num{3.3e9}$ वर्ष, 3.3 अरब वर्ष। (3) यह पृथ्वी की आयु से कम है,
इसलिए चाप पूरी \omterm{def:g12:continental-crust:continental-crust}{महाद्वीपीय भूपर्पटी} बना सकते थे; पर
\omterm{def:g11:subduction-zones:subduction}{क्षेपण} भूपर्पटी को वापस भी ले जाता है, इसलिए वास्तव में भूपर्पटी का बहुत-सा भाग एक से अधिक बार
बना और नष्ट हुआ है।
\end{solution}

\begin{solution}{exo:g12:subduction-magmatism:11}
(1) उथले भूकंप ऊपर चढ़ने वाली प्लेट में और दोनों प्लेटों के बीच के भ्रंश पर
भी होते हैं; केवल गहरे भूकंप अकेले \omterm{def:g11:subduction-zones:slab}{स्लैब} में पड़ते हैं।
(2) \qty{86.4}{km} पर $116 - 80 = \qty{36}{km}$: $\tan\theta = 36/86.4
\approx 0.417$, $\theta \approx \ang{22.6}$।
\end{solution}

\begin{solution}{exo:g12:subduction-magmatism:12}
प्रस्ताव के समकोण त्रिभुज में \omterm{def:g11:subduction-zones:slab}{स्लैब} कर्ण $L$ है
और $d$ कोण $\theta$ से लगी भुजा, इसलिए $\cos\theta = d/L$ और $L =
d/\cos\theta$। जापान के लिए $L = 270/\cos\ang{20} \approx 270/0.940 \approx
\qty{287}{km}$। प्रति वर्ष \qty{83}{mm} पर: $\num{2.87e8}/83 \approx \num{3.5e6}$
वर्ष, लगभग 3.5 मिलियन वर्ष।
\end{solution}

\begin{solution}{exo:g12:subduction-magmatism:13}
(1) $V = \pi r^2 h = \pi \times 0.5^2 \times 8 \approx \qty{6.3}{km^3}$।
(2) हिमस्खलन के आयतन का लगभग $6.3/2.5 \approx 2.5$ गुना।
(3) $P = \rho g z$: \qty{6}{km} पर $2830 \times 9.81 \times 6000 \approx
\qty{1.67e8}{Pa} = \qty{0.17}{GPa}$; \qty{14}{km} पर लगभग
\qty{3.89e8}{Pa} = \qty{0.39}{GPa}।
\end{solution}

\begin{solution}{exo:g12:subduction-magmatism:14}
(1) $2.7 - 3.2 = \qty{-0.5}{km^3}$ प्रति वर्ष, प्रति मिलियन वर्ष \qty{5e5}{km^3}
की हानि। (2) $500 \times \num{5e5} = \qty{2.5e8}{km^3}$, यानी
भूपर्पटी का $\num{2.5e8}/\num{9e9} \approx \qty{2.8}{\percent}$। (3) हर
दर एक अनुमान है, जिसकी अनिश्चितता उनके बीच के प्रति वर्ष \qty{0.5}{km^3} के
अंतर से कहीं अधिक है: त्रुटियों के भीतर दोनों संतुलित हैं।
\end{solution}

\begin{solution}{exo:g12:subduction-magmatism:15}
(1) $\num{20000} \times 10 = \qty{2e5}{km^3}$। (2) संसार के सभी महाद्वीपीय चापों को मिलाकर $\num{2e5}/1 =
\num{200000}$ वर्ष। किसी एक
चाप को उस मैग्मा का केवल छोटा-सा भाग मिलता है, इसलिए एक \omterm{def:g12:subduction-magmatism:batholith}{बैथोलिथ} को
कहीं अधिक समय लगता है: वह \omterm{def:g11:igneous-rocks:intrusion}{प्लूटॉन} दर \omterm{def:g11:igneous-rocks:intrusion}{प्लूटॉन}, कई मिलियन वर्षों में बनता है। (3) कम से
कम उसके शीर्ष के ऊपर की \qty{5}{km} चट्टान।
\end{solution}

\begin{solution}{pb:g12:subduction-magmatism:1}
\textbf{1.} $111.2 \times \cos\ang{39} \approx 111.2 \times 0.777 \approx
\qty{86.4}{km}$।

\textbf{2.} $(144.1 - 141.0) \times 86.4 \approx 3.1 \times 86.4 \approx
\qty{268}{km}$, लगभग \qty{270}{km}।

\textbf{3.} $(141.0 - 139.7) \times 86.4 \approx \qty{112}{km}$।

\textbf{4.} ऊपर चढ़ने वाली प्लेट पर, जो नहीं डूबती, खाई के
पश्चिम में: प्रशांत प्लेट उसके नीचे डूबती है।

\textbf{5.} वहाँ \omterm{def:g11:subduction-zones:slab}{स्लैब} अभी इतनी गहरी नहीं है कि अपना जल
पिघलने लायक गरम प्रावार में छोड़ सके।

\textbf{6.} $116 - 80 = \qty{36}{km}$।

\textbf{7.} देशांतर का एक अंश, \qty{86.4}{km}।

\textbf{8.} $\tan\theta = 36/86.4 \approx 0.417$, इसलिए $\theta \approx
\ang{22.6}$।

\textbf{9.} रेखा पूर्व की ओर हर \qty{86.4}{km} पर \qty{36}{km} उठती है, इसलिए वह
इवाते से $98/0.417 \approx \qty{235}{km}$ पूर्व में सतह तक पहुँचती है,
$141.0 + 235/86.4 \approx \ang{143.7}$~पू पर: खाई से लगभग $(144.1 - 143.7) \times 86.4
\approx \qty{35}{km}$ पश्चिम में।

\textbf{10.} \omterm{def:g11:subduction-zones:slab}{स्लैब} एक समतल तल नहीं है: वह खाई में लगभग
क्षैतिज प्रवेश करती है और डूबते हुए अधिकाधिक नीचे मुड़ती जाती है।

\textbf{11.} $268 \times \tan\ang{20} \approx 268 \times 0.364 \approx
\qty{98}{km}$।

\textbf{12.} वे सहमत हैं: दोनों तरीक़ों से \qty{98}{km}।

\textbf{13.} दोनों संसार के परास के भीतर हैं, \qty{105}{km} के माध्य से
थोड़े उथले।

\textbf{14.} मेगाता इवाते से \qty{112}{km} पश्चिम में है, इसलिए उसके नीचे \omterm{def:g11:subduction-zones:slab}{स्लैब}
लगभग $98 + 0.417 \times 112 \approx \qty{145}{km}$ गहरी है।

\textbf{15.} प्रावार में, भूपर्पटी के आधार से लगभग $98 - 35 = \qty{63}{km}$
नीचे।

\textbf{16.} उसने उस कटक के निकट प्रवेश किया जहाँ भूपर्पटी बनी, समुद्री
जल से; उसका बहुत-सा भाग सर्पेंटाइन में बंद है।

\textbf{17.} $T = -5.104 \times 3.16^2 + 132.9 \times 3.16 + 1120.7 \approx
\qty{1490}{\celsius}$।

\textbf{18.} $1490 - 1000 = 490$ से $1490 - 950 = \qty{540}{\celsius}$ तक:
लगभग \qty{500}{\celsius}।

\textbf{19.} \omterm{def:g11:subduction-zones:slab}{स्लैब} अपने चारों ओर के प्रावार से कई सौ अंश ठंडी है, जल के साथ भी
पिघलने के लिए बहुत ठंडी: उसके ऊपर की गरम प्रावार-कील,
भीगने पर, पिघलती है।

\textbf{20.} उत्तर-पूर्वी जापान की \textbf{ज्वालामुखीय अग्र-रेखा के नीचे \omterm{def:g11:subduction-zones:slab}{स्लैब} लगभग \qty{100}{km}
पर है} (दोनों विधियों से \qty{98}{km})।
\end{solution}
